\begin{lemma}
For any fixed red row map $\rho: \operatorname{Fin}n \to \operatorname{Fin}m$, let
\[
X_{\rm all}
=\{\phi:\operatorname{Fin}n\hookrightarrow
\operatorname{Fin}m\times\operatorname{Fin}m:
(\phi(v))_1=\rho(v)\text{ for every }v\}.
\]
For each row $i\in\operatorname{Fin}m$, put $R_i=\{v:\rho(v)=i\}$ and $q_i=|R_i|$.
Then the cardinality of $X_{\rm all}$ is given by the product over rows of $(m)_{q_i}$, where $(m)_{q_i}$ is the falling factorial:
\[
|X_{\rm all}| = \prod_{i\in\operatorname{Fin}m} (m)_{q_i}.
\]
\end{lemma}
\begin{proof}
For each row \(i\), let \(\mathcal B_i\) be the set of injective maps
\(b_i:R_i\hookrightarrow \operatorname{Fin}m\).  We first give an explicit
bijection
\[
X_{\rm all}\longleftrightarrow \prod_{i\in\operatorname{Fin}m}\mathcal B_i .
\]
Given \(\phi\in X_{\rm all}\), define its row restriction \(b_i\) by
\[
b_i(v)=(\phi(v))_2\qquad (v\in R_i).
\]
This is injective on \(R_i\): if \(v,w\in R_i\) and \(b_i(v)=b_i(w)\), then
\((\phi(v))_2=(\phi(w))_2\), while the fixed-row condition gives
\((\phi(v))_1=\rho(v)=i=\rho(w)=(\phi(w))_1\).  Hence
\(\phi(v)=\phi(w)\), and injectivity of \(\phi\) gives \(v=w\).

Conversely, suppose we are given a family of rowwise injections
\((b_i)_i\), with \(b_i:R_i\hookrightarrow\operatorname{Fin}m\).  Define
\[
\phi(v)=(\rho(v), b_{\rho(v)}(v)).
\]
Here \(v\in R_{\rho(v)}\) by definition of \(R_{\rho(v)}\), so the expression
is well-defined.  The first coordinate of \(\phi(v)\) is \(\rho(v)\), so
\(\phi\) satisfies the fixed-row condition.  To prove \(\phi\) is globally
injective, suppose \(\phi(v)=\phi(w)\).  Equality of first coordinates gives
\(\rho(v)=\rho(w)=i\).  Thus \(v,w\in R_i\).  Equality of second coordinates
then gives \(b_i(v)=b_i(w)\), and injectivity of \(b_i\) gives \(v=w\).
These two constructions are inverse to each other, since restricting the
assembled map recovers each \(b_i\), and assembling the row restrictions of
\(\phi\) recovers \(\phi\) pointwise.

It remains to count each row factor.  If \(q_i>m\), then there is no
injective map from the \(q_i\)-element set \(R_i\) into
\(\operatorname{Fin}m\), so \(|\mathcal B_i|=0\).  The falling factorial
\((m)_{q_i}\) is also \(0\) in this case.  If \(q_i\le m\), then an element
of \(\mathcal B_i\) is counted by choosing distinct blue coordinates for the
vertices of \(R_i\), one at a time: there are \(m\) choices for the first
vertex, \(m-1\) for the second, and so on.  Thus
\[
(m)_{q_i}=m(m-1)\cdots(m-q_i+1),
\]
choices are available in row \(i\).  Therefore \(|\mathcal B_i|=(m)_{q_i}\)
for every row, including the zero case \(q_i>m\).  The displayed bijection
and the finite product rule give
\[
|X_{\rm all}|=\prod_{i\in\operatorname{Fin}m}|\mathcal B_i|
=\prod_{i\in\operatorname{Fin}m}(m)_{q_i},
\]
as claimed.
\end{proof}
